% --- Archon physics formalization source begin ---
% archon:physics
% archon:covers PhyXMiniProblems/problem_phyx_mini_0584.lean
% archon:source-report reports/phyx_mini/problem_phyx_mini_0584.source.json

\chapter{Physics problem phyx\_mini\_0584}
\label{ch:PhyXMiniProblems_problem_phyx_mini_0584}

\paragraph{Problem source.}
\#\# Physical scenario

Figure shows a two-dimensional, infinite-potential well lying in an \textbackslash{}( xy \textbackslash{}) plane that contains an electron. We probe for the electron along a line that bisects \textbackslash{}( L\_x \textbackslash{}) and find three points at which the detection probability is maximum. Those points are separated by \textbackslash{}( 2.00\textbackslash{}text\{nm\} \textbackslash{}). Then we probe along a line that bisects \textbackslash{}( L\_y \textbackslash{}) and find five points at which the detection probability is maximum. Those points are separated by \textbackslash{}( 3.00\textbackslash{}text\{nm\} \textbackslash{}).

\#\# Question

What is the energy of the electron?

\#\# Answer choices

- A: A: \textbackslash{}frac\{n\textasciicircum{}2 h\textasciicircum{}2\}\{2 \textbackslash{}pi\textasciicircum{}2 m d\textasciicircum{}2\}

- B: B: \textbackslash{}frac\{n\textasciicircum{}2 h\textasciicircum{}2\}\{4 \textbackslash{}pi\textasciicircum{}2 m d\textasciicircum{}2\}

- C: C: \textbackslash{}frac\{n\textasciicircum{}2 h\textasciicircum{}2\}\{8 \textbackslash{}pi\textasciicircum{}2 m d\textasciicircum{}2\}

- D: D: \textbackslash{}frac\{n h\textasciicircum{}2\}\{4 \textbackslash{}pi\textasciicircum{}2 m d\textasciicircum{}2\}

\#\# Figure caption (auxiliary; use the image as primary evidence)

The image depicts a rectangular box positioned in a two-dimensional coordinate plane with axes labeled \textbackslash{}(x\textbackslash{}) and \textbackslash{}(y\textbackslash{}).

1. **Rectangle Attributes**:
   - The rectangle's edges are emphasized with a thick blue outline.
   - The width along the x-axis is indicated by a dashed line marked \textbackslash{}(L\_x\textbackslash{}).
   - The height along the y-axis is shown with a dashed line labeled \textbackslash{}(L\_y\textbackslash{}).

2. **Axes**:
   - The y-axis is vertical on the left.
   - The x-axis is horizontal at the bottom.

3. **Lines**:
   - A vertical dashed line runs from the top to the bottom of the rectangle.
   - A horizontal dashed line extends from the left to the right side of the rectangle.

4. **Text Labels**:
   - \textbackslash{}(L\_x\textbackslash{}): Positioned horizontally near the top border of the rectangle.
   - \textbackslash{}(L\_y\textbackslash{}): Positioned vertically near the right border.

The overall configuration presents a clear and simple geometric layout with emphasis on dimensions \textbackslash{}(L\_x\textbackslash{}) and \textbackslash{}(L\_y\textbackslash{}).

\#\# Dataset metadata

Quantum Phenomena; Physical Model Grounding Reasoning, Multi-Formula Reasoning

\paragraph{Recorded answer/context.}
C: C: \textbackslash{}frac\{n\textasciicircum{}2 h\textasciicircum{}2\}\{8 \textbackslash{}pi\textasciicircum{}2 m d\textasciicircum{}2\}

\paragraph{Figure/image path.}
/root/proposal\_for\_physic/science-mango/phyx\_mini\_run/phyx\_data/test\_image/584.png

\paragraph{Formalization target.}
create a compiling Lean file with sorry bodies at `PhyXMiniProblems/problem\_phyx\_mini\_0584.lean`. The Lean declarations must preserve the physical quantities, dimensions or dimensional roles, figure labels, governing-law hypotheses, and final relation expressed by this problem.
Use Mathlib/Physlib names found through LeanExplore where available. If a physics API is missing, introduce faithful local abstractions rather than scalar placeholder aliases.

\begin{theorem}[Physics formalization target]
\label{thm:physics:phyx_mini_0584:target}
\lean{PhyXMiniProblems.ProblemPhyXMini0584.problem_phyx_mini_0584}
\uses{def:physics:phyx-mini-0584:phyxminiproblems-problemphyxmini0584-energyinjoules, def:physics:phyx-mini-0584:phyxminiproblems-problemphyxmini0584-energyinelectronvolts, def:physics:phyx-mini-0584:phyxminiproblems-problemphyxmini0584-rectangularwellelectronsetup, def:physics:phyx-mini-0584:phyxminiproblems-problemphyxmini0584-matchesrectangularwellelectronscenario, def:physics:phyx-mini-0584:phyxminiproblems-problemphyxmini0584-matchesprimaryrectangularwellfigure, def:physics:phyx-mini-0584:phyxminiproblems-problemphyxmini0584-matchesdetectionprobereadouts, def:physics:phyx-mini-0584:phyxminiproblems-problemphyxmini0584-useselectronreferencedata, def:physics:phyx-mini-0584:phyxminiproblems-problemphyxmini0584-usesstandardplanckconstant, def:physics:phyx-mini-0584:phyxminiproblems-problemphyxmini0584-hasphysicalrectangularwellparameters, def:physics:phyx-mini-0584:phyxminiproblems-problemphyxmini0584-satisfiesinfiniterectangularwelllaws}
The assigned autoformalize agent should translate this physics problem into Lean declarations in the covered file, with theorem and lemma proof bodies written as `by sorry`.
\end{theorem}
\begin{proof}
This is an autoformalization task, not a proof task. Produce statements that can later be proved by the physics prover without weakening the physical model.
\end{proof}

% --- Archon declaration topology begin ---
\section{Declaration topology}
\begin{definition}[Length Quantity]
  \label{def:physics:phyx-mini-0584:phyxminiproblems-problemphyxmini0584-lengthquantity}
  \lean{PhyXMiniProblems.ProblemPhyXMini0584.LengthQuantity}
  A nonnegative, unit-independent physical length
\end{definition}
\begin{proof}
  This is the declared model definition of Length Quantity.
\end{proof}
\begin{definition}[Mass Quantity]
  \label{def:physics:phyx-mini-0584:phyxminiproblems-problemphyxmini0584-massquantity}
  \lean{PhyXMiniProblems.ProblemPhyXMini0584.MassQuantity}
  A nonnegative, unit-independent physical mass
\end{definition}
\begin{proof}
  This is the declared model definition of Mass Quantity.
\end{proof}
\begin{definition}[action Dimension]
  \label{def:physics:phyx-mini-0584:phyxminiproblems-problemphyxmini0584-actiondimension}
  \lean{PhyXMiniProblems.ProblemPhyXMini0584.actionDimension}
  The physical dimension of action, equivalently energy times time
\end{definition}
\begin{proof}
  This is the declared model definition of action Dimension.
\end{proof}
\begin{definition}[Action Quantity]
  \label{def:physics:phyx-mini-0584:phyxminiproblems-problemphyxmini0584-actionquantity}
  \lean{PhyXMiniProblems.ProblemPhyXMini0584.ActionQuantity}
  \uses{def:physics:phyx-mini-0584:phyxminiproblems-problemphyxmini0584-actiondimension}
  A nonnegative, unit-independent physical action such as Planck's constant
\end{definition}
\begin{proof}
  This is the declared model definition of Action Quantity.
\end{proof}
\begin{definition}[length Readout]
  \label{def:physics:phyx-mini-0584:phyxminiproblems-problemphyxmini0584-lengthreadout}
  \lean{PhyXMiniProblems.ProblemPhyXMini0584.lengthReadout}
  \uses{def:physics:phyx-mini-0584:phyxminiproblems-problemphyxmini0584-lengthquantity}
  Read a physical length in a selected length unit
\end{definition}
\begin{proof}
  This is the declared model definition of length Readout.
\end{proof}
\begin{definition}[length In Meters]
  \label{def:physics:phyx-mini-0584:phyxminiproblems-problemphyxmini0584-lengthinmeters}
  \lean{PhyXMiniProblems.ProblemPhyXMini0584.lengthInMeters}
  \uses{def:physics:phyx-mini-0584:phyxminiproblems-problemphyxmini0584-lengthquantity, def:physics:phyx-mini-0584:phyxminiproblems-problemphyxmini0584-lengthreadout}
  Metre readout of a physical length
\end{definition}
\begin{proof}
  This is the declared model definition of length In Meters.
\end{proof}
\begin{definition}[length In Nanometers]
  \label{def:physics:phyx-mini-0584:phyxminiproblems-problemphyxmini0584-lengthinnanometers}
  \lean{PhyXMiniProblems.ProblemPhyXMini0584.lengthInNanometers}
  \uses{def:physics:phyx-mini-0584:phyxminiproblems-problemphyxmini0584-lengthquantity, def:physics:phyx-mini-0584:phyxminiproblems-problemphyxmini0584-lengthreadout}
  Nanometre readout used for the supplied probe separations
\end{definition}
\begin{proof}
  This is the declared model definition of length In Nanometers.
\end{proof}
\begin{definition}[mass In Kilograms]
  \label{def:physics:phyx-mini-0584:phyxminiproblems-problemphyxmini0584-massinkilograms}
  \lean{PhyXMiniProblems.ProblemPhyXMini0584.massInKilograms}
  \uses{def:physics:phyx-mini-0584:phyxminiproblems-problemphyxmini0584-massquantity}
  Kilogram readout of a physical mass
\end{definition}
\begin{proof}
  This is the declared model definition of mass In Kilograms.
\end{proof}
\begin{definition}[action In Joule Seconds]
  \label{def:physics:phyx-mini-0584:phyxminiproblems-problemphyxmini0584-actioninjouleseconds}
  \lean{PhyXMiniProblems.ProblemPhyXMini0584.actionInJouleSeconds}
  \uses{def:physics:phyx-mini-0584:phyxminiproblems-problemphyxmini0584-actionquantity}
  Joule-second readout of a physical action
\end{definition}
\begin{proof}
  This is the declared model definition of action In Joule Seconds.
\end{proof}
\begin{definition}[energy In Joules]
  \label{def:physics:phyx-mini-0584:phyxminiproblems-problemphyxmini0584-energyinjoules}
  \lean{PhyXMiniProblems.ProblemPhyXMini0584.energyInJoules}
  Joule readout of a physical energy
\end{definition}
\begin{proof}
  This is the declared model definition of energy In Joules.
\end{proof}
\begin{definition}[energy In Electron Volts]
  \label{def:physics:phyx-mini-0584:phyxminiproblems-problemphyxmini0584-energyinelectronvolts}
  \lean{PhyXMiniProblems.ProblemPhyXMini0584.energyInElectronVolts}
  \uses{def:physics:phyx-mini-0584:phyxminiproblems-problemphyxmini0584-energyinjoules}
  Electron-volt readout grounded in Physlib's calibrated electron volt
\end{definition}
\begin{proof}
  This is the declared model definition of energy In Electron Volts.
\end{proof}
\begin{definition}[Coordinate Axis]
  \label{def:physics:phyx-mini-0584:phyxminiproblems-problemphyxmini0584-coordinateaxis}
  \lean{PhyXMiniProblems.ProblemPhyXMini0584.CoordinateAxis}
  The two coordinate axes of the rectangular well
\end{definition}
\begin{proof}
  This is the declared model definition of Coordinate Axis.
\end{proof}
\begin{definition}[to Fin]
  \label{def:physics:phyx-mini-0584:phyxminiproblems-problemphyxmini0584-coordinateaxis-tofin}
  \lean{PhyXMiniProblems.ProblemPhyXMini0584.CoordinateAxis.toFin}
  \uses{def:physics:phyx-mini-0584:phyxminiproblems-problemphyxmini0584-coordinateaxis}
  The Fin 2 coordinate selected by an axis label
\end{definition}
\begin{proof}
  This is the declared model definition of to Fin.
\end{proof}
\begin{definition}[coordinate In Meters]
  \label{def:physics:phyx-mini-0584:phyxminiproblems-problemphyxmini0584-coordinateinmeters}
  \lean{PhyXMiniProblems.ProblemPhyXMini0584.coordinateInMeters}
  \uses{def:physics:phyx-mini-0584:phyxminiproblems-problemphyxmini0584-coordinateaxis}
  Read one coordinate of a point in the coherent metre chart
\end{definition}
\begin{proof}
  This is the declared model definition of coordinate In Meters.
\end{proof}
\begin{definition}[Particle Species]
  \label{def:physics:phyx-mini-0584:phyxminiproblems-problemphyxmini0584-particlespecies}
  \lean{PhyXMiniProblems.ProblemPhyXMini0584.ParticleSpecies}
  Particle species confined by the potential
\end{definition}
\begin{proof}
  This is the declared model definition of Particle Species.
\end{proof}
\begin{definition}[Spatial Plane]
  \label{def:physics:phyx-mini-0584:phyxminiproblems-problemphyxmini0584-spatialplane}
  \lean{PhyXMiniProblems.ProblemPhyXMini0584.SpatialPlane}
  Spatial plane in which the well lies
\end{definition}
\begin{proof}
  This is the declared model definition of Spatial Plane.
\end{proof}
\begin{definition}[Confinement Model]
  \label{def:physics:phyx-mini-0584:phyxminiproblems-problemphyxmini0584-confinementmodel}
  \lean{PhyXMiniProblems.ProblemPhyXMini0584.ConfinementModel}
  Confinement idealization stated in the problem
\end{definition}
\begin{proof}
  This is the declared model definition of Confinement Model.
\end{proof}
\begin{definition}[Probe Line]
  \label{def:physics:phyx-mini-0584:phyxminiproblems-problemphyxmini0584-probeline}
  \lean{PhyXMiniProblems.ProblemPhyXMini0584.ProbeLine}
  The two dashed probe lines. verticalBisectorOfLx fixes x = L\_x/2 and varies y; horizontalBisectorOfLy fixes y = L\_y/2 and varies x
\end{definition}
\begin{proof}
  This is the declared model definition of Probe Line.
\end{proof}
\begin{definition}[varied Axis]
  \label{def:physics:phyx-mini-0584:phyxminiproblems-problemphyxmini0584-probeline-variedaxis}
  \lean{PhyXMiniProblems.ProblemPhyXMini0584.ProbeLine.variedAxis}
  \uses{def:physics:phyx-mini-0584:phyxminiproblems-problemphyxmini0584-coordinateaxis, def:physics:phyx-mini-0584:phyxminiproblems-problemphyxmini0584-probeline}
  Coordinate that varies while moving along a probe line
\end{definition}
\begin{proof}
  This is the declared model definition of varied Axis.
\end{proof}
\begin{definition}[bisected Axis]
  \label{def:physics:phyx-mini-0584:phyxminiproblems-problemphyxmini0584-probeline-bisectedaxis}
  \lean{PhyXMiniProblems.ProblemPhyXMini0584.ProbeLine.bisectedAxis}
  \uses{def:physics:phyx-mini-0584:phyxminiproblems-problemphyxmini0584-coordinateaxis, def:physics:phyx-mini-0584:phyxminiproblems-problemphyxmini0584-probeline}
  Coordinate fixed at half the corresponding side length
\end{definition}
\begin{proof}
  This is the declared model definition of bisected Axis.
\end{proof}
\begin{definition}[Figure Label]
  \label{def:physics:phyx-mini-0584:phyxminiproblems-problemphyxmini0584-figurelabel}
  \lean{PhyXMiniProblems.ProblemPhyXMini0584.FigureLabel}
  Literal mathematical labels visible in the supplied raster
\end{definition}
\begin{proof}
  This is the declared model definition of Figure Label.
\end{proof}
\begin{definition}[Rectangular Well Figure]
  \label{def:physics:phyx-mini-0584:phyxminiproblems-problemphyxmini0584-rectangularwellfigure}
  \lean{PhyXMiniProblems.ProblemPhyXMini0584.RectangularWellFigure}
  \uses{def:physics:phyx-mini-0584:phyxminiproblems-problemphyxmini0584-figurelabel}
  Primary-image features of the rectangular well and its two bisectors
\end{definition}
\begin{proof}
  This is the declared model definition of Rectangular Well Figure.
\end{proof}
\begin{definition}[Detection Probe Observation]
  \label{def:physics:phyx-mini-0584:phyxminiproblems-problemphyxmini0584-detectionprobeobservation}
  \lean{PhyXMiniProblems.ProblemPhyXMini0584.DetectionProbeObservation}
  \uses{def:physics:phyx-mini-0584:phyxminiproblems-problemphyxmini0584-lengthquantity}
  A finite readout of all equal-height detection-probability maxima found on one probe line. The points use the coherent metre coordinate chart; their common spacing remains a dimensionful length
\end{definition}
\begin{proof}
  This is the declared model definition of Detection Probe Observation.
\end{proof}
\begin{definition}[Rectangular Well Electron Setup]
  \label{def:physics:phyx-mini-0584:phyxminiproblems-problemphyxmini0584-rectangularwellelectronsetup}
  \lean{PhyXMiniProblems.ProblemPhyXMini0584.RectangularWellElectronSetup}
  \uses{def:physics:phyx-mini-0584:phyxminiproblems-problemphyxmini0584-lengthquantity, def:physics:phyx-mini-0584:phyxminiproblems-problemphyxmini0584-massquantity, def:physics:phyx-mini-0584:phyxminiproblems-problemphyxmini0584-actionquantity, def:physics:phyx-mini-0584:phyxminiproblems-problemphyxmini0584-coordinateaxis, def:physics:phyx-mini-0584:phyxminiproblems-problemphyxmini0584-particlespecies, def:physics:phyx-mini-0584:phyxminiproblems-problemphyxmini0584-spatialplane, def:physics:phyx-mini-0584:phyxminiproblems-problemphyxmini0584-confinementmodel, def:physics:phyx-mini-0584:phyxminiproblems-problemphyxmini0584-probeline, def:physics:phyx-mini-0584:phyxminiproblems-problemphyxmini0584-rectangularwellfigure, def:physics:phyx-mini-0584:phyxminiproblems-problemphyxmini0584-detectionprobeobservation}
  Independent physical quantities for the confined electron. In particular, stateEnergy, quantum numbers, and side lengths are fields constrained by general laws below; none is defined from the requested answer
\end{definition}
\begin{proof}
  This is the declared model definition of Rectangular Well Electron Setup.
\end{proof}
\begin{definition}[Is Inside Rectangular Well]
  \label{def:physics:phyx-mini-0584:phyxminiproblems-problemphyxmini0584-isinsiderectangularwell}
  \lean{PhyXMiniProblems.ProblemPhyXMini0584.IsInsideRectangularWell}
  \uses{def:physics:phyx-mini-0584:phyxminiproblems-problemphyxmini0584-lengthinmeters, def:physics:phyx-mini-0584:phyxminiproblems-problemphyxmini0584-coordinateinmeters, def:physics:phyx-mini-0584:phyxminiproblems-problemphyxmini0584-rectangularwellelectronsetup}
  A point lies in the closed rectangular well in the SI coordinate chart
\end{definition}
\begin{proof}
  This is the declared model definition of Is Inside Rectangular Well.
\end{proof}
\begin{definition}[Lies On Probe Line]
  \label{def:physics:phyx-mini-0584:phyxminiproblems-problemphyxmini0584-liesonprobeline}
  \lean{PhyXMiniProblems.ProblemPhyXMini0584.LiesOnProbeLine}
  \uses{def:physics:phyx-mini-0584:phyxminiproblems-problemphyxmini0584-lengthinmeters, def:physics:phyx-mini-0584:phyxminiproblems-problemphyxmini0584-coordinateinmeters, def:physics:phyx-mini-0584:phyxminiproblems-problemphyxmini0584-probeline, def:physics:phyx-mini-0584:phyxminiproblems-problemphyxmini0584-rectangularwellelectronsetup}
  A point lies on the selected geometric bisector of the well
\end{definition}
\begin{proof}
  This is the declared model definition of Lies On Probe Line.
\end{proof}
\begin{definition}[detection Probability Density]
  \label{def:physics:phyx-mini-0584:phyxminiproblems-problemphyxmini0584-detectionprobabilitydensity}
  \lean{PhyXMiniProblems.ProblemPhyXMini0584.detectionProbabilityDensity}
  \uses{def:physics:phyx-mini-0584:phyxminiproblems-problemphyxmini0584-rectangularwellelectronsetup}
  Pointwise detection-probability density of the chosen wave representative
\end{definition}
\begin{proof}
  This is the declared model definition of detection Probability Density.
\end{proof}
\begin{definition}[Is Detection Probability Maximum Along]
  \label{def:physics:phyx-mini-0584:phyxminiproblems-problemphyxmini0584-isdetectionprobabilitymaximumalong}
  \lean{PhyXMiniProblems.ProblemPhyXMini0584.IsDetectionProbabilityMaximumAlong}
  \uses{def:physics:phyx-mini-0584:phyxminiproblems-problemphyxmini0584-probeline, def:physics:phyx-mini-0584:phyxminiproblems-problemphyxmini0584-rectangularwellelectronsetup, def:physics:phyx-mini-0584:phyxminiproblems-problemphyxmini0584-isinsiderectangularwell, def:physics:phyx-mini-0584:phyxminiproblems-problemphyxmini0584-liesonprobeline, def:physics:phyx-mini-0584:phyxminiproblems-problemphyxmini0584-detectionprobabilitydensity}
  A global detection-probability maximum along one in-well probe segment
\end{definition}
\begin{proof}
  This is the declared model definition of Is Detection Probability Maximum Along.
\end{proof}
\begin{definition}[Matches Rectangular Well Electron Scenario]
  \label{def:physics:phyx-mini-0584:phyxminiproblems-problemphyxmini0584-matchesrectangularwellelectronscenario}
  \lean{PhyXMiniProblems.ProblemPhyXMini0584.MatchesRectangularWellElectronScenario}
  \uses{def:physics:phyx-mini-0584:phyxminiproblems-problemphyxmini0584-rectangularwellelectronsetup}
  Categorical facts explicitly stated by the physical scenario
\end{definition}
\begin{proof}
  This is the declared model definition of Matches Rectangular Well Electron Scenario.
\end{proof}
\begin{definition}[Matches Primary Rectangular Well Figure]
  \label{def:physics:phyx-mini-0584:phyxminiproblems-problemphyxmini0584-matchesprimaryrectangularwellfigure}
  \lean{PhyXMiniProblems.ProblemPhyXMini0584.MatchesPrimaryRectangularWellFigure}
  \uses{def:physics:phyx-mini-0584:phyxminiproblems-problemphyxmini0584-rectangularwellfigure}
  Geometry and labels transcribed from the primary raster
\end{definition}
\begin{proof}
  This is the declared model definition of Matches Primary Rectangular Well Figure.
\end{proof}
\begin{definition}[Matches Detection Probe Readouts]
  \label{def:physics:phyx-mini-0584:phyxminiproblems-problemphyxmini0584-matchesdetectionprobereadouts}
  \lean{PhyXMiniProblems.ProblemPhyXMini0584.MatchesDetectionProbeReadouts}
  \uses{def:physics:phyx-mini-0584:phyxminiproblems-problemphyxmini0584-lengthinmeters, def:physics:phyx-mini-0584:phyxminiproblems-problemphyxmini0584-lengthinnanometers, def:physics:phyx-mini-0584:phyxminiproblems-problemphyxmini0584-coordinateinmeters, def:physics:phyx-mini-0584:phyxminiproblems-problemphyxmini0584-probeline, def:physics:phyx-mini-0584:phyxminiproblems-problemphyxmini0584-rectangularwellelectronsetup, def:physics:phyx-mini-0584:phyxminiproblems-problemphyxmini0584-detectionprobabilitydensity, def:physics:phyx-mini-0584:phyxminiproblems-problemphyxmini0584-isdetectionprobabilitymaximumalong}
  The two experimental scan readouts. Besides the stated counts and spacings, the fields say that the listed points really are distinct maxima, exhaust all maxima on the relevant line, occur with nonzero density, and are listed in increasing order with the reported uniform adjacent separation
\end{definition}
\begin{proof}
  This is the declared model definition of Matches Detection Probe Readouts.
\end{proof}
\begin{definition}[Uses Electron Reference Data]
  \label{def:physics:phyx-mini-0584:phyxminiproblems-problemphyxmini0584-useselectronreferencedata}
  \lean{PhyXMiniProblems.ProblemPhyXMini0584.UsesElectronReferenceData}
  \uses{def:physics:phyx-mini-0584:phyxminiproblems-problemphyxmini0584-massinkilograms, def:physics:phyx-mini-0584:phyxminiproblems-problemphyxmini0584-rectangularwellelectronsetup}
  Reference mass needed for numerical evaluation. No electron-mass constant was found in the searched Physlib API, so the 2022 CODATA kilogram readout is stated separately from the problem and figure observations
\end{definition}
\begin{proof}
  This is the declared model definition of Uses Electron Reference Data.
\end{proof}
\begin{definition}[Uses Standard Planck Constant]
  \label{def:physics:phyx-mini-0584:phyxminiproblems-problemphyxmini0584-usesstandardplanckconstant}
  \lean{PhyXMiniProblems.ProblemPhyXMini0584.UsesStandardPlanckConstant}
  \uses{def:physics:phyx-mini-0584:phyxminiproblems-problemphyxmini0584-actioninjouleseconds, def:physics:phyx-mini-0584:phyxminiproblems-problemphyxmini0584-rectangularwellelectronsetup}
  Physlib supplies the reduced Planck constant Constants.ℏ in joule-seconds; ordinary Planck action is h = 2πℏ
\end{definition}
\begin{proof}
  This is the declared model definition of Uses Standard Planck Constant.
\end{proof}
\begin{definition}[Has Physical Rectangular Well Parameters]
  \label{def:physics:phyx-mini-0584:phyxminiproblems-problemphyxmini0584-hasphysicalrectangularwellparameters}
  \lean{PhyXMiniProblems.ProblemPhyXMini0584.HasPhysicalRectangularWellParameters}
  \uses{def:physics:phyx-mini-0584:phyxminiproblems-problemphyxmini0584-lengthinmeters, def:physics:phyx-mini-0584:phyxminiproblems-problemphyxmini0584-massinkilograms, def:physics:phyx-mini-0584:phyxminiproblems-problemphyxmini0584-actioninjouleseconds, def:physics:phyx-mini-0584:phyxminiproblems-problemphyxmini0584-energyinjoules, def:physics:phyx-mini-0584:phyxminiproblems-problemphyxmini0584-coordinateaxis, def:physics:phyx-mini-0584:phyxminiproblems-problemphyxmini0584-probeline, def:physics:phyx-mini-0584:phyxminiproblems-problemphyxmini0584-rectangularwellelectronsetup}
  Positivity and nondegeneracy conditions for the physical setup
\end{definition}
\begin{proof}
  This is the declared model definition of Has Physical Rectangular Well Parameters.
\end{proof}
\begin{definition}[Satisfies Infinite Rectangular Well Laws]
  \label{def:physics:phyx-mini-0584:phyxminiproblems-problemphyxmini0584-satisfiesinfiniterectangularwelllaws}
  \lean{PhyXMiniProblems.ProblemPhyXMini0584.SatisfiesInfiniteRectangularWellLaws}
  \uses{def:physics:phyx-mini-0584:phyxminiproblems-problemphyxmini0584-lengthinmeters, def:physics:phyx-mini-0584:phyxminiproblems-problemphyxmini0584-massinkilograms, def:physics:phyx-mini-0584:phyxminiproblems-problemphyxmini0584-actioninjouleseconds, def:physics:phyx-mini-0584:phyxminiproblems-problemphyxmini0584-energyinjoules, def:physics:phyx-mini-0584:phyxminiproblems-problemphyxmini0584-coordinateinmeters, def:physics:phyx-mini-0584:phyxminiproblems-problemphyxmini0584-probeline, def:physics:phyx-mini-0584:phyxminiproblems-problemphyxmini0584-rectangularwellelectronsetup, def:physics:phyx-mini-0584:phyxminiproblems-problemphyxmini0584-isinsiderectangularwell}
  The general stationary-state laws for a two-dimensional infinite rectangular well. The pointwise representative has the separable sine profile inside the hard walls and vanishes outside. Along a non-nodal bisector scan, the number of equal maxima is the quantum number for the varying coordinate, adjacent maxima are separated by L/n, and the spectrum is E\_(n\_x,n\_y) = h\textasciicircum{}2/(8m) * (n\_x\textasciicircum{}2/L\_x\textasciicircum{}2 + n\_y\textasciicircum{}2/L\_y\textasciicircum{}2). These are general relations involving the independent setup fields;
\end{definition}
\begin{proof}
  This is the declared model definition of Satisfies Infinite Rectangular Well Laws.
\end{proof}
\begin{definition}[Answer Choice]
  \label{def:physics:phyx-mini-0584:phyxminiproblems-problemphyxmini0584-answerchoice}
  \lean{PhyXMiniProblems.ProblemPhyXMini0584.AnswerChoice}
  The four answer labels printed in the supplied text
\end{definition}
\begin{proof}
  This is the declared model definition of Answer Choice.
\end{proof}
\begin{definition}[displayed Symbolic Energy Formula]
  \label{def:physics:phyx-mini-0584:phyxminiproblems-problemphyxmini0584-displayedsymbolicenergyformula}
  \lean{PhyXMiniProblems.ProblemPhyXMini0584.displayedSymbolicEnergyFormula}
  \uses{def:physics:phyx-mini-0584:phyxminiproblems-problemphyxmini0584-answerchoice}
  Literal symbolic expression printed beside each answer. The variables are dimensionless mode number n, action readout h, mass readout m, and length readout d; this definition records the source verbatim and is not a physical law used by the theorem
\end{definition}
\begin{proof}
  This is the declared model definition of displayed Symbolic Energy Formula.
\end{proof}
\begin{definition}[recorded Dataset Answer]
  \label{def:physics:phyx-mini-0584:phyxminiproblems-problemphyxmini0584-recordeddatasetanswer}
  \lean{PhyXMiniProblems.ProblemPhyXMini0584.recordedDatasetAnswer}
  \uses{def:physics:phyx-mini-0584:phyxminiproblems-problemphyxmini0584-answerchoice}
  Answer label recorded by the dataset, retained independently of the model
\end{definition}
\begin{proof}
  This is the declared model definition of recorded Dataset Answer.
\end{proof}
\begin{lemma}[quantum Numbers from probe peak counts]
  \label{lem:physics:phyx-mini-0584:phyxminiproblems-problemphyxmini0584-quantumnumbers-from-probe-peak-counts}
  \lean{PhyXMiniProblems.ProblemPhyXMini0584.quantumNumbers_from_probe_peak_counts}
  \uses{def:physics:phyx-mini-0584:phyxminiproblems-problemphyxmini0584-rectangularwellelectronsetup, def:physics:phyx-mini-0584:phyxminiproblems-problemphyxmini0584-matchesdetectionprobereadouts, def:physics:phyx-mini-0584:phyxminiproblems-problemphyxmini0584-satisfiesinfiniterectangularwelllaws}
  The five maxima on the horizontal line determine n\_x = 5, while the three maxima on the vertical line determine n\_y = 3
\end{lemma}
\begin{proof}
  The conclusion follows from the governing relations and figure readouts named in the statement of quantum Numbers from probe peak counts.
\end{proof}
\begin{lemma}[side Lengths from probe peak spacings]
  \label{lem:physics:phyx-mini-0584:phyxminiproblems-problemphyxmini0584-sidelengths-from-probe-peak-spacings}
  \lean{PhyXMiniProblems.ProblemPhyXMini0584.sideLengths_from_probe_peak_spacings}
  \uses{def:physics:phyx-mini-0584:phyxminiproblems-problemphyxmini0584-lengthinnanometers, def:physics:phyx-mini-0584:phyxminiproblems-problemphyxmini0584-rectangularwellelectronsetup, def:physics:phyx-mini-0584:phyxminiproblems-problemphyxmini0584-matchesdetectionprobereadouts, def:physics:phyx-mini-0584:phyxminiproblems-problemphyxmini0584-satisfiesinfiniterectangularwelllaws}
  Since adjacent maxima of sin\textasciicircum{}2(nπx/L) are separated by L/n, the scans give L\_x = 5(3 nm) = 15 nm and L\_y = 3(2 nm) = 6 nm
\end{lemma}
\begin{proof}
  The conclusion follows from the governing relations and figure readouts named in the statement of side Lengths from probe peak spacings.
\end{proof}
\begin{lemma}[energy from probe spacings]
  \label{lem:physics:phyx-mini-0584:phyxminiproblems-problemphyxmini0584-energy-from-probe-spacings}
  \lean{PhyXMiniProblems.ProblemPhyXMini0584.energy_from_probe_spacings}
  \uses{def:physics:phyx-mini-0584:phyxminiproblems-problemphyxmini0584-massinkilograms, def:physics:phyx-mini-0584:phyxminiproblems-problemphyxmini0584-actioninjouleseconds, def:physics:phyx-mini-0584:phyxminiproblems-problemphyxmini0584-energyinjoules, def:physics:phyx-mini-0584:phyxminiproblems-problemphyxmini0584-rectangularwellelectronsetup, def:physics:phyx-mini-0584:phyxminiproblems-problemphyxmini0584-matchesdetectionprobereadouts, def:physics:phyx-mini-0584:phyxminiproblems-problemphyxmini0584-hasphysicalrectangularwellparameters, def:physics:phyx-mini-0584:phyxminiproblems-problemphyxmini0584-satisfiesinfiniterectangularwelllaws}
  Substituting the two scan spacings into the two-dimensional spectrum eliminates the mode numbers and side lengths: n\_x/L\_x = 1/(3 nm) and n\_y/L\_y = 1/(2 nm)
\end{lemma}
\begin{proof}
  The conclusion follows from the governing relations and figure readouts named in the statement of energy from probe spacings.
\end{proof}
% --- Archon declaration topology end ---
% --- Archon physics formalization source end ---
